\documentclass{report}
\usepackage[utf8]{inputenc}
\usepackage[T1]{fontenc}
\usepackage{geometry}
\usepackage{graphicx}
\usepackage{float}
\usepackage{hyperref}
\usepackage[english,indonesian]{babel}

\title{Optimasi Pemrosesan Template LaTeX Menggunakan Pendekatan Berbasis AI}
\author{Academic Support Agent}
\date{\today}

\begin{document}
\maketitle

\begin{abstract}

\end{abstract}

\chapter{PENDAHULUAN}


\chapter{TINJAUAN PUSTAKA}


% Original content (commented out):
% \documentclass[12pt,graphicx,caption,rotating,letterpaper,oneside]{report}
\usepackage[utf8]{inputenc}
\usepackage[activeacute,spanish]{babel}
\usepackage{graphicx}
\usepackage[dvipsnames,usenames]{color}
\usepackage{colortbl}
\usepackage{subfigure}
\usepackage[pdftex,colorlinks=true,linkcolor=black,urlcolor=black]{hyperref}
\usepackage{url}
\usepackage{indentfirst}
\usepackage[T1]{fontenc}
\usepackage{ulem}
%\usepackage{efont,mathesf}
\usepackage{tgpagella}  %Bonita Palatino enriquecida pero no con formulas
\setlength{\voffset}{-0.5cm}
\setlength{\hoffset}{0.9cm}
\setlength{\headsep}{0pt}
\setlength{\headheight}{0pt}
\setlength{\oddsidemargin}{-0.8in}
\setlength{\marginparwidth}{-0.5cm}
\setlength{\textwidth}{18cm}
\setlength{\footskip}{2pt}
\setlength{\topmargin}{0in}
\setlength{\textheight}{27cm}
\setlength{\fboxrule}{3pt}
\usepackage[T1]{fontenc}
\usepackage[sc]{mathpazo}

\begin{document}
\thispagestyle{empty}
\begin{tabular}{p{2.4cm}p{15cm}}
\includegraphics[width=2.4cm]{uadylogo.png}
\begin{center}
\rule[2cm]{0.8mm}{19.5cm}%vertical
\hspace{1pt}
\rule[2cm]{0.4mm}{19.5cm}%vertical
\end{center}
&
\vspace{-3.5cm}
\begin{center}
\rule[1mm]{15cm}{0.4mm}%horizontal
\vspace{0.1pt}
\rule[3mm]{15cm}{0.8mm}%horizontal
\\
\LARGE{\bf{Universidad Autonóma de Yucatán}}\\
\vspace*{0.3cm}
\large{\bf{Facultad de Ingeniería}}

\vspace{4.5\baselineskip}
{\Large \bf{Caracterización térmica de  }}\\
\vspace*{0.2cm}
{\Large \bf{ferrofluidos con cargas de polvo de diamante }}\\
\vspace*{0.2cm}
\Large \bf{ bajo la influencia de campos magnéticos.}\\
\vspace*{0.2cm}
\vspace*{1.6cm}
\Large{\bf TESIS}

\vspace*{1.5cm}
{\large \bf{Presentada por:}}

\vspace*{0.7cm}
{\LARGE \bf{Graciela Inés Quiñones Weiss}}


\vspace*{2cm}
{\large \bf{En su examen profesional}}\\
\vspace*{0.2cm}
{\large \bf{En opción al título de:}}\\
\vspace*{0.5cm}
\huge{\bf Ingeniero Físico}

\vspace*{1.8cm}

{\large \bf{Mérida, Yucatán, México}}\\
\vspace*{-0.4cm}
{\large \bf{2013}}
\end{center}
\end{tabular}

% \newpage
\thispagestyle{empty}
$ $
\end{document} 

\end{document}